% Generated draft: requires mathematical and source review.
\documentclass[11pt]{article}
\usepackage[margin=25mm]{geometry}
\usepackage{fontspec}
\setmainfont{Latin Modern Roman}
\usepackage{amsmath,amssymb,mathtools,unicode-math}
\setmathfont{Latin Modern Math}
\usepackage{enumitem,needspace,xurl,hyperref}
\setlength{\parindent}{0pt}
\setlength{\parskip}{5pt}
\setlength{\emergencystretch}{3em}
\newcommand{\N}{\mathbb{N}}
\newcommand{\Z}{\mathbb{Z}}
\newcommand{\Q}{\mathbb{Q}}
\newcommand{\R}{\mathbb{R}}
\newcommand{\C}{\mathbb{C}}
\newcommand{\F}{\mathbb{F}}
\newcommand{\A}{\mathbb{A}}
\newcommand{\PP}{\mathbb{P}}
\newcommand{\E}{\mathbb{E}}
\newcommand{\eps}{\varepsilon}
\newcommand{\dd}{\,\mathrm{d}}
\begin{document}
\section*{UnsolvedMath ID: 30005920}
\textbf{Source problem code:} catalogue entry
\subsection{Definitions and mathematical statement}
Under the conventional complex-matrix interpretation, let \(n,m\geq 1\) be integers, let \(M_n=M_n(\C)\), and write \(M_n^{\mathrm{sa}}=\{X\in M_n:X^*=X\}\). Let
\[
A=(A_1,\ldots,A_m),\qquad B=(B_1,\ldots,B_m)\in(M_n^{\mathrm{sa}})^m.
\]
The scalar field and the ranges of \(n,m\) are implicit in the supplied statement; the conventions above make its usual matrix interpretation explicit.

A complex-linear map \(\Phi:M_n\to M_n\) is completely positive if, for every integer \(r\geq 1\), the map \(\operatorname{id}_{M_r}\otimes\Phi\) sends positive semidefinite elements of \(M_r\otimes M_n\) to positive semidefinite elements. It is unital if \(\Phi(I_n)=I_n\), and trace preserving if \(\operatorname{Tr}(\Phi(X))=\operatorname{Tr}(X)\) for every \(X\in M_n\), where \(\operatorname{Tr}\) denotes the ordinary matrix trace.

The supplied interpolation question asks for necessary and sufficient conditions on \(A\) and \(B\) for the existence of a unital quantum channel \(\Phi:M_n\to M_n\) satisfying
\[
\Phi(B_k)=A_k\qquad\text{for every }k\in\{1,\ldots,m\}.
\]
The term ``unital quantum channel'' is not defined in the supplied excerpt. Under the convention that a quantum channel is a completely positive, trace-preserving map, the question is to characterize when there exists a completely positive complex-linear map satisfying
\[
\Phi(I_n)=I_n,\qquad
\operatorname{Tr}(\Phi(X))=\operatorname{Tr}(X)\quad(X\in M_n),\qquad
\Phi(B_k)=A_k\quad(1\leq k\leq m).
\]
If instead the report uses ``quantum channel'' in the observable-picture sense of a unital completely positive map, the trace-preservation condition is absent. These are distinct formulations; the supplied evidence does not settle which convention the report intends. No commutativity or other structural hypothesis on either tuple is stated.

The intended unresolved target also requires clarification: the supplied verification note says that the report follows this question with a necessary-and-sufficient characterization in Theorem~5, whereas the supplied dated assessment labels the problem partially solved and says a general classification for arbitrary self-adjoint tuples is not known. The statement and scope of Theorem~5 are not supplied. Consequently, this draft records the interpolation question without asserting that its existence characterization remains open or specifying an unsupported further classification problem.
\subsection{Short English statement}
Given two equally sized tuples of self-adjoint matrices, characterize when one unital quantum channel sends every \(B_k\) to the corresponding \(A_k\). The source excerpt leaves the channel convention unclear, particularly whether trace preservation is required. A supplied note reports an existing necessary-and-sufficient theorem, so the precise remaining open target needs review.
\subsection{Sources}
\begin{itemize}
\item Supplied local UnsolvedMath registry record, permanent ID \(30005920\), exact source problem code \texttt{catalogue entry}, titled ``Quantum-Channel Interpolation Characterization.'' The record belongs to the Oberwolfach Reports Open Problems collection and attributes the problem to the 2024 report; no individual proposer is supplied.
\item \emph{Non-commutative Function Theory and Free Probability} (2024), Oberwolfach Reports, pp.~1227--1284, DOI \texttt{10.4171/owr/2024/22}: \url{https://doi.org/10.4171/owr/2024/22}. The supplied problem excerpt carries the page marker 1257. The background describes these reports as documenting mathematical workshops and their open-problem discussions. This report citation is supplied provenance and has not been independently consulted here.
\item Li--Poon: \url{https://arxiv.org/abs/1012.1675}.
\item Li--Poon full paper: \url{https://people.wm.edu/~cklixx/cpafinal.pdf}.
\item Literature assessment dated exactly \texttt{2026-08-22}: supplied dated curation derived from the OWR record, not independently verified. It labels the problem ``partially solved'' and reports Choi-matrix and operator-system criteria, substantial results for structured or commuting families, and no known general classification for arbitrary self-adjoint tuples. Its time-sensitive status and cited claims require independent verification.
\item The supplied statement-verification note says that the official report poses the interpolation question and then gives a necessary-and-sufficient characterization as Theorem~5. The registry nevertheless marks \texttt{statement\_status} as \texttt{unrecoverable}; no theorem text or explanation of the remaining target is included.
\end{itemize}
\end{document}
